\section{函数的最大值最小值问题}\label{sec:5.4}

我们先来看一个具体实例．

一块边长为 $a$ 的正方形，在四个角上截去同样大小的正方形，做成无盖的盒，问截去多大的平方块能使盒的容积最大？

诸如此类求函数最大值或最小值的问题，在现实生活中随处可见．本节我们就来讨论这类问题．

设函数 $f(x)$ 在 $(a, b)$ 内的 $x_0$ 点达到最大（小）值，由最大（小）值的定义知 $x_0$ 必是 $f(x)$ 在 $(a, b)$ 的极大（小）值点，从而 $x_0$ 是 $f(x)$ 的稳定点或不可导点．此外，函数的最大（小）值也可能在区间的端点达到．因此可能的最大（小）值点是 $f(x)$ 的稳定点，不可导点或区间的端点．

\begin{figure}[htbp]
\centering
\includegraphics[width=0.28\textwidth]{fig5-13.png}
\figcap{图 5-13}\label{fig:5-13}
\end{figure}

但是，函数的极值是函数在某点的局部性态，对同一函数，极小值甚至可能大于极大值，因此极大值未必是函数在某区间上的最大（小）值．下面就两种常见的情形给出判别法，以最大值为例说明．

\textbf{1. 闭区间情形}\quad 设 $f(x)$ 在 $[a, b]$ 连续，这时 $f(x)$ 在 $[a, b]$ 必有最大值．若稳定点和不可导点只有有限个，则将所有稳定点、不可导点和区间端点的函数值进行比较，最大者即是最大值．

\textbf{2. 开区间情形}\quad 设 $f(x)$ 在 $(a, b)$ 可导，且在 $(a, b)$ 有最大值．若在 $(a, b)$ 内有唯一的稳定点 $x_0$，则 $x_0$ 是最大值点．

下面我们来求解上面提出的容积最大的无盖盒问题．

设 $x$ 为截去的小方块的边长，则盒的容积为
\[
V(x) = x(a - 2x)^2,\quad x \in \left(0, \frac{a}{2}\right).
\]
显然，$V(x)$ 在 $\left(0, \frac{a}{2}\right)$ 可导，且
\[
V'(x) = (a - 2x)^2 - 4x(a - 2x) = (a - 2x)(a - 6x).
\]
令 $V'(x) = 0$ 得 $x = \frac{a}{2}$ 或 $x = \frac{a}{6}$．因此在 $\left(0, \frac{a}{2}\right)$ 中有唯一的稳定点 $\frac{a}{6}$．由实际问题本身知 $V(x)$ 在 $\left(0, \frac{a}{2}\right)$ 中必有最大值，故知最大值为 $V\left(\frac{a}{6}\right) = \frac{2}{27}a^3$．即截去的小的方块边长为 $\frac{a}{6}$ 时，盒的容积最大．

\begin{example}
求函数 $f(x) = \left|2x^3 - 9x^2 + 12x\right|$ 在 $[-1, 3]$ 的最大值和最小值．
\end{example}

\begin{proof}[解]
\[
2x^3 - 9x^2 + 12x = x\left[2\left(x - \frac{9}{4}\right)^2 + \frac{15}{8}\right],
\]
因此
\[
f(x) = (2x^3 - 9x^2 + 12x)\operatorname{sgn} x,\quad x \in [-1, 3],
\]
\[
f'(x) = (6x^2 - 18x + 12)\operatorname{sgn} x = 6(x - 1)(x - 2)\operatorname{sgn} x,\quad x \in (-1, 0) \cup (0, 3),
\]
故 $f(x)$ 的稳定点为 $x = 1$，$x = 2$，不可导点为 $x = 0$．
\[
f(-1) = 23,\quad f(0) = 0,\quad f(1) = 5,\quad f(2) = 4,\quad f(3) = 9.
\]
最大值为 $f(-1) = 23$，最小值为 $f(0) = 0$．
\end{proof}

\begin{example}
做一个圆柱形无盖铁桶，容积一定，设为 $V_0$．问铁桶的底半径与高的比例应为多少，才能最省铁皮？
\end{example}

\begin{proof}[解]
设铁桶底半径为 $r$，高为 $h$（见图~\ref{fig:5-14}），则所需铁皮面积为
\[
s = 2\pi r h + \pi r^2.
\]
利用已知条件 $V_0 = \pi r^2 h$，得 $h = \frac{V_0}{\pi r^2}$．则面积 $s$ 可化为 $r$ 的函数
\[
s(r) = \frac{2V_0}{r} + \pi r^2,\quad 0 < r < +\infty.
\]
于是问题化为求函数 $s$ 在 $(0, +\infty)$ 内的最小值问题．
\[
s'(r) = -\frac{2V_0}{r^2} + 2\pi r = \frac{2\pi r^3 - 2V_0}{r^2}.
\]
令 $s'(r) = 0$，得到唯一的稳定点 $r_0 = \sqrt[3]{\dfrac{V_0}{\pi}}$，又由实际问题本身知 $s(r)$ 在 $(0, +\infty)$ 必有最小值，从而唯一的稳定点 $r_0$ 必是最小值点，此时有
\[
h = \left.\frac{V_0}{\pi r^2}\right|_{r = r_0} = \sqrt[3]{\frac{V_0}{\pi}} = r_0,
\]
即当底半径 $r$ 与高 $h$ 相等，均为 $\sqrt[3]{\dfrac{V_0}{\pi}}$ 时，最省铁皮．
\end{proof}

\begin{figure}[htbp]
\centering
\includegraphics[width=0.3\textwidth]{fig5-14.png}
\figcap{图 5-14}\label{fig:5-14}
\end{figure}

\begin{example}
在正午时，甲船恰在乙船正南 $82\ \mathrm{km}$ 处，以速度 $v_1 = 20\ \mathrm{km/h}$ 向正东开出；乙船也正以速度 $v_2 = 16\ \mathrm{km/h}$ 向正南开去（图 5-15）．已知两船航向不变，试证：下午二时，两船相距最近．
\end{example}

\begin{figure}[htbp]
\centering
\includegraphics[width=0.45\textwidth]{fig5-15.png}
\figcap{图 5-15}\label{fig:5-15}
\end{figure}

\begin{proof}[证明]
设 $t$ 小时后，两船相距 $y\ \mathrm{km}$，则显然有
\[
y(t) = \sqrt{(20t)^2 + (82 - 16t)^2} = 2\sqrt{164t^2 - 656t + 1681},
\]
\[
y' = \frac{328t - 656}{\sqrt{164t^2 - 656t + 1681}},\quad 0 \leqslant t < +\infty.
\]
由于 $y \neq 0$，因此 $y'$ 的分母不为 $0$，即函数在 $(0, +\infty)$ 可导．令 $y' = 0$，得唯一稳定点 $t = 2$，比较 $t = 0$ 和 $t = 2$ 点的值，
\[
y(0) = 82,\quad y(2) = \sqrt{41} \times 10,
\]
显然
\[
\lim_{t \to +\infty} y(t) = +\infty,
\]
故 $t = 2$ 时函数达到最小值，即下午二时，两船相距最近．
\end{proof}

\begin{example}
根据物理学的费马原理，光线沿着所需时间为最少的路线传播．今有 I，II 两种介质，以 $L$ 为分界线．光在介质 I 与介质 II 中的传播速度分别为 $v_1$ 与 $v_2$．问：光线由介质 I 中的点 $A$ 到介质 II 中的点 $B$，应走哪一条路线？
\end{example}

\begin{proof}[解]
取分界线 $L$ 所在直线为 $Ox$ 轴．过 $A$，$B$ 作 $L$ 的垂线，设垂足为 $A_1$，$B_1$，并设 $AA_1 = a$，$BB_1 = b$，$A_1B_1 = c$，选定 $A_1$ 为坐标原点 $O$（图~\ref{fig:5-16}）．

\begin{figure}[htbp]
\centering
\includegraphics[width=0.45\textwidth]{fig5-16.png}
\figcap{图 5-16}\label{fig:5-16}
\end{figure}

光线在同一介质中的传播途径应当是直线．设想光线从点 $A$ 到点 $B$ 所走的路线通过 $L$ 上的点 $M$，$M$ 的坐标为 $x$．于是问题化为，当 $x$ 取何值时，折线 $AMB$ 才是光线所走的路线．

光线从点 $A$ 到达点 $B$ 所需的时间为
\[
t = \frac{AM}{v_1} + \frac{BM}{v_2}
= \frac{\sqrt{a^2 + x^2}}{v_1} + \frac{\sqrt{b^2 + (c - x)^2}}{v_2}
\quad (-\infty < x < +\infty).
\]
根据费马原理，我们要求的是上述函数 $t = f(x)$ 的最小值．
\[
\frac{\mathrm{d}t}{\mathrm{d}x} = f'(x)
= \frac{x}{v_1\sqrt{a^2 + x^2}} - \frac{c - x}{v_2\sqrt{b^2 + (c - x)^2}},
\]
\[
\frac{\mathrm{d}^2 t}{\mathrm{d}x^2} = f''(x)
= \frac{a^2}{v_1\left(a^2 + x^2\right)^{3/2}}
+ \frac{b^2}{v_2\left[b^2 + (c - x)^2\right]^{3/2}}.
\]
因为 $f''(x)$ 恒为正，所以 $f'(x)$ 在 $(-\infty, +\infty)$ 上严格单调上升，从而方程 $f'(x) = 0$ 至多有一个根，即函数 $t = f(x)$ 至多有一个稳定点．又因为 $f'(x)$ 是 $x$ 的连续函数，且
\[
f'(0) = \frac{-c}{v_2\sqrt{b^2 + c^2}} < 0,\qquad
f'(c) = \frac{c}{v_1\sqrt{a^2 + c^2}} > 0,
\]
所以方程 $f'(x) = 0$ 的根位于区间 $(0, c)$ 内，记作 $x_0$．这就是函数 $t = f(x)$ 的唯一稳定点．已知 $f''(x)$ 恒为正，因此 $f''(x_0) > 0$，于是由极值第二充分条件，$f(x_0)$ 为函数 $t = f(x)$ 的极小值．又
\[
\lim_{x \to +\infty} f(x) = +\infty,\qquad \lim_{x \to -\infty} f(x) = +\infty.
\]
因而连续函数 $t = f(x)$ 的最小值必在 $(-\infty, +\infty)$ 内部达到．于是可以断定，唯一的极小值 $f(x_0)$ 就是最小值．这表明，当点 $M$ 的横坐标 $x = x_0$ 时，折线 $AMB$ 就是光线所走的路线．

上面的讨论只告诉我们：$x_0 \in (0, c)$，并不知道 $x_0$ 的具体数值．求出 $x_0$ 的值比较困难，不过实际上并不需要，我们可以从几何上作如下说明．

$x_0$ 所满足的方程
\[
f'(x) = \frac{x}{v_1\sqrt{a^2 + x^2}} - \frac{c - x}{v_2\sqrt{b^2 + (c - x)^2}} = 0
\]
可写为
\[
\frac{1}{v_1}\cdot\frac{A_1M}{AM} = \frac{1}{v_2}\cdot\frac{B_1M}{BM},
\]
即
\[
\frac{\sin\alpha}{v_1} = \frac{\sin\beta}{v_2},
\]
或
\[
\frac{\sin\alpha}{\sin\beta} = \frac{v_1}{v_2}.
\]
这就是说，入射角与折线角的正弦之比等于光在两介质中的传播速度之比，这正是光学上的折射定律．上面的讨论说明，光在不同的两种介质中传播时，遵守折射定律．
\end{proof}

\begin{xiti}
\item 求下列函数在指定区间上的最大值与最小值：
\begin{ziti}
\item $y = x^5 - 5x^4 + 5x^3 + 1$，$[-1, 2]$；
\item $y = 2\tan x - \tan^2 x$，$[0, \frac{\pi}{2})$；
\item $y = \sqrt{x}\ln x$，$(0, +\infty)$；
\item $y = |x^2 - 3x + 2|$，$[-10, 10]$；
\item $y = e^{|x-3|}$，$[-5, 5]$．
\end{ziti}
\item 给定长为 $l$ 的线段，试把它分成两段，使以这两段为边所围成的矩形面积为最大．
\item 设用某仪器进行测量时，读得 $n$ 次实验数据为 $a_1, a_2, \dots, a_n$，问以怎样的数值 $x$ 表达所要测量的真值，才能使它与这 $n$ 个数之差的平方和为最小？
\item 求内接于椭圆 $\frac{x^2}{a^2} + \frac{y^2}{b^2} = 1$ 而边平行于坐标轴的面积最大的矩形．
\item 求点 $M(p, p)$ 到抛物线 $y^2 = 2px$ 的最短距离．
\item 做一个圆柱形锅炉，已知其容积为 $V$，两端面材料的每单位面积价格为 $a$ 元，侧面材料的每单位面积价格为 $b$ 元，问锅炉的直径与高的比等于多少时，造价最省？
\item 某村计划修建一条断面面积为 $4\ \mathrm{m}^2$ 的梯形渠道，侧面的坡度为 $\frac{3}{4}$（即底边与斜高间夹角 $\theta$ 满足 $\tan\theta = \frac{3}{4}$），底边 $b$ 与斜高 $l$ 为多长时湿周最小．（根据经验，湿周最小时渠道过水能力最大．）
\item 设炮口的仰角为 $\alpha$，炮弹的初速为 $v_0\ \mathrm{m/s}$，炮口取作原点，发炮时间取作 $t = 0$，不计空气阻力时，炮弹的运动方程为：
\[
\begin{cases}
x = t v_0 \cos\alpha, \\
y = t v_0 \sin\alpha - \frac{1}{2}gt^2.
\end{cases}
\]
若初速 $v_0$ 不变，问如何调整炮口的仰角 $\alpha$，使炮弹射程最远？
\end{xiti}
